% =========================================================================
% DOCUMENTCLASS
% =========================================================================
\documentclass[12pt,oneside]{book}

% =========================================================================
% METADATOS
% =========================================================================
\newcommand{\universidad}{Universidad de Buenos Aires (UBA)}
\newcommand{\facultad}{Facultad de Derecho}
\newcommand{\carrera}{Carrera de Abogacía}
\newcommand{\catedra}{Cátedra de Práctica Profesional}
\newcommand{\materia}{Derecho Procesal Civil y Comercial}
\newcommand{\titulo}{La demanda: estructura, competencia y organización judicial}
\newcommand{\autor}{Isaac Edgar Camacho Ocampo}
\newcommand{\anio}{2026}

% =========================================================================
% PAQUETES
% =========================================================================
\usepackage[utf8]{inputenc}
\usepackage[T1]{fontenc}
\usepackage[spanish,es-nodecimaldot]{babel}

\usepackage[
    top=2.5cm,
    bottom=2.5cm,
    left=3cm,
    right=2.5cm,
    headheight=15pt
]{geometry}

\usepackage{amsmath}
\usepackage{amssymb}
\usepackage{mathtools}

\usepackage{graphicx}
\usepackage{float}
\usepackage{caption}
\usepackage{subcaption}

\usepackage{booktabs}
\usepackage{array}
\usepackage{longtable}
\usepackage{tabularx}

\usepackage{enumitem}

\usepackage{xcolor}
\usepackage[most]{tcolorbox}

\usepackage{fancyhdr}

\usepackage{ragged2e}

% =========================================================================
% CONFIGURACIÓN
% =========================================================================
\graphicspath{
    {./img/}
    {./graficos/}
    {./figuras/}
}

\setcounter{tocdepth}{2}

\setlength{\parindent}{0pt}
\setlength{\parskip}{0.5em}

\setlist[itemize]{
    topsep=0.3em,
    itemsep=0.2em
}

\setlist[enumerate]{
    topsep=0.3em,
    itemsep=0.2em
}

\clubpenalty=10000
\widowpenalty=10000

% =========================================================================
% COMANDOS
% =========================================================================
\newcommand{\abs}[1]{\left|#1\right|}
\newcommand{\norm}[1]{\left\lVert#1\right\rVert}
\newcommand{\vect}[1]{\mathbf{#1}}
\newcommand{\deriv}[2]{\frac{d#1}{d#2}}
\newcommand{\pderiv}[2]{\frac{\partial#1}{\partial#2}}

% =========================================================================
% ENTORNOS / CAJAS
% =========================================================================
\newtcolorbox{definition}[1]{
    enhanced,
    breakable,
    title={Definición: #1},
    fonttitle=\bfseries,
    colback=blue!3,
    colframe=blue!50!black,
    boxrule=0.8pt,
    arc=2mm
}

\newtcolorbox{important}{
    enhanced,
    breakable,
    title={Importante},
    fonttitle=\bfseries,
    colback=yellow!8,
    colframe=orange!70!black,
    boxrule=0.8pt,
    arc=2mm
}

\newtcolorbox{example}[1]{
    enhanced,
    breakable,
    title={Ejemplo: #1},
    fonttitle=\bfseries,
    colback=green!4,
    colframe=green!40!black,
    boxrule=0.8pt,
    arc=2mm
}

\newtcolorbox{formula}[1]{
    enhanced,
    breakable,
    title={Fórmula / Regla: #1},
    fonttitle=\bfseries,
    colback=purple!4,
    colframe=purple!50!black,
    boxrule=0.8pt,
    arc=2mm
}

\newtcolorbox{exercise}[1]{
    enhanced,
    breakable,
    title={Ejercicio: #1},
    fonttitle=\bfseries,
    colback=gray!5,
    colframe=gray!60!black,
    boxrule=0.8pt,
    arc=2mm
}

\newtcolorbox{note}{
    enhanced,
    breakable,
    title={Nota},
    fonttitle=\bfseries,
    colback=white,
    colframe=gray!50,
    boxrule=0.6pt,
    arc=2mm
}

% =========================================================================
% CONFIGURACIÓN FINAL
% =========================================================================
\pagestyle{fancy}
\fancyhf{}

\fancyhead[L]{\nouppercase{\leftmark}}
\fancyhead[R]{\materia}

\fancyfoot[C]{\thepage}

\renewcommand{\headrulewidth}{0.4pt}
\renewcommand{\footrulewidth}{0pt}

\fancypagestyle{plain}{
    \fancyhf{}
    \fancyfoot[C]{\thepage}
    \renewcommand{\headrulewidth}{0pt}
}

\usepackage[
    colorlinks=true,
    linkcolor=blue!50!black,
    citecolor=green!40!black,
    urlcolor=blue!60!black,
    pdftitle={\titulo},
    pdfauthor={\autor},
    pdfsubject={Apuntes universitarios},
    bookmarks=true
]{hyperref}

% =========================================================================
% DOCUMENT
% =========================================================================
\begin{document}

% -------------------------------------------------------------------------
% PORTADA
% -------------------------------------------------------------------------
\begin{titlepage}

\thispagestyle{empty}

\begin{center}

\vspace*{1cm}

{\large \textsc{\universidad}}

\vspace{0.2cm}

{\large \textsc{\facultad}}

\vspace{0.2cm}

{\large \carrera}

\vspace{0.2cm}

{\normalsize \catedra}

\vspace{3cm}

{\Huge \textbf{\materia}}

\vspace{1cm}

{\LARGE \titulo}

\vspace{0.8cm}

\textit{Comprender la lógica detrás de cada palabra de un escrito judicial es el primer paso para ejercer el derecho con criterio propio.}

\vspace{1.2cm}

\rule{0.70\textwidth}{0.5pt}

\vspace{0.5cm}

{\large Apuntes de estudio}

\vspace{0.5cm}

\rule{0.70\textwidth}{0.5pt}

\vfill

{\large \autor}

\vspace{0.3cm}

{\large \anio}

\end{center}

\vfill

\begin{flushleft}
\textit{Este es un modesto aporte a los estudiantes de «Derecho Procesal Civil y Comercial» no pretende sustituir el material oficial de la materia.}
\end{flushleft}

\end{titlepage}

% -------------------------------------------------------------------------
% ÍNDICE
% -------------------------------------------------------------------------
\tableofcontents

% -------------------------------------------------------------------------
% PRESENTACIÓN DE LA UNIDAD
% -------------------------------------------------------------------------
\chapter*{Presentación de la unidad}
\addcontentsline{toc}{chapter}{Presentación de la unidad}

La presente unidad aborda la construcción de una demanda civil y comercial a partir de sus fundamentos teóricos, su estructura formal y las decisiones de competencia que su presentación exige. El desarrollo recorre, de manera integrada, los siguientes temas:

\begin{itemize}
    \item Jurisdicción, proceso y el rol del abogado litigante.
    \item El tema decidendum y la relación entre el hecho y la norma.
    \item Retórica jurídica y preguntas retóricas.
    \item Estructura de la demanda conforme al artículo 330 del CPCCN.
    \item Legitimación activa y pasiva.
    \item Redacción de los hechos.
    \item Competencia territorial, federal y ordinaria.
    \item La organización histórica de la justicia en la Ciudad Autónoma de Buenos Aires y el fallo «Vinas».
\end{itemize}

\begin{note}
Todo abogado litigante deberá redactar una demanda desde el comienzo de su ejercicio profesional. El valor de esta unidad no reside en proporcionar un formulario para copiar, sino en explicar la lógica que subyace a cada una de las partes de la demanda: por qué se narra primero el hecho y después se pide el derecho, por qué el objeto limita lo que el juez puede conceder, y a quién y ante qué fuero conviene demandar. Comprender la diferencia entre el «ser» y el «deber ser», dominar la subsunción y utilizar la retórica de manera honesta constituyen herramientas de cada escrito, cada audiencia y cada negociación profesional, más allá de la especialidad que se elija. Distinguir un hecho de una interpretación, fundamentar un reclamo con claridad y advertir la arbitrariedad de una decisión ajena son, asimismo, competencias que exceden el ámbito judicial y resultan aplicables a cualquier conflicto cotidiano.
\end{note}

A lo largo de la unidad se trabaja, a modo de caso guía, un accidente ocurrido entre una lancha comercial de pasajeros y una pequeña embarcación en el río Tigre, que sirve de base para ejemplificar la construcción práctica de una demanda de daños y perjuicios.

% =========================================================================
% CAPÍTULO 1
% =========================================================================
\chapter[Fundamentos del proceso judicial]{Fundamentos del proceso judicial}

\section{Jurisdicción: decir y aplicar el derecho}

La jurisdicción constituye la potestad que tiene el juez y que no tiene ningún otro funcionario: consiste en \textbf{decir} el derecho aplicable a un caso y \textbf{aplicarlo}. Por esa razón existe la sentencia: si la parte condenada por ella no cumple espontáneamente con la consecuencia allí establecida, el juez puede hacérsela cumplir, incluso por la fuerza.

\begin{definition}{Jurisdicción}
Potestad exclusiva del juez de decir el derecho aplicable a un caso y de aplicarlo, incluso haciéndolo cumplir por la fuerza si la parte condenada no lo hace espontáneamente.
\end{definition}

\section{El objetivo del proceso y la garantía de igualdad}

El proceso constituye una estructura particular orientada a un objetivo concreto: arribar a una sentencia que resuelva el conflicto planteado entre las partes. A través de ese mecanismo, el proceso procura garantizar la \textbf{igualdad} entre las partes, en tanto la justicia, como ideal absoluto, resulta \textit{«un ideal discutible, inalcanzable, quizás imposible de acceder para el ser humano»}. Lo que puede proponerse, con humildad, es un procedimiento lo más equilibrado posible, que coloque a ambas partes frente a un tercero imparcial encargado de dirimir el conflicto.

\section{El rol del abogado litigante frente al rol del juez}

El rol de los magistrados es relevante, pero el de los operadores jurídicos —quienes efectivamente reclaman y ponen en tela de juicio una cuestión— resulta tan importante, o a veces más importante, que el del magistrado. El fallo judicial es la consecuencia de un litigio previamente planteado por las partes: se llega a él \textbf{litigando}, entendiendo el litigio como el conflicto y la manera en que las partes lo plantean.

\section{El tema decidendum: los límites de la decisión judicial}

Entre la demanda y la contestación existe un límite que condiciona la decisión judicial.

\begin{definition}{Tema decidendum}
Límite dentro del cual puede resolver el juez, fijado por la pretensión de la parte actora, la defensa del demandado y la prueba producida en el proceso.
\end{definition}

De allí la utilidad de estudiar el planteo que realizan los abogados —y no solamente la sentencia final—, ya que sin un verdadero desarrollo previo no puede alcanzarse ese final; de lo contrario, la sentencia sería, según se señala en el apunte, \textit{«algo así como el final de la película sin haber visto todo el desarrollo»}. El entrenamiento de todo operador jurídico debe apuntar a que ese desarrollo previo sea lo más trascendente, eficaz y útil posible.

\section{El mundo del ser y el mundo del deber ser}

Frente a un hecho, es necesario pensarlo y encuadrarlo en categorías jurídicas para que resulte idóneo a los fines del derecho. La filosofía del derecho distingue, a este respecto, dos mundos ontológicos:

\begin{itemize}
    \item El \textbf{mundo del ser}: integrado por los hechos, es decir, por los acontecimientos de la realidad. Por ejemplo, que una persona tome un arma o cause un daño.
    \item El \textbf{mundo del deber ser}: integrado por las normas, es decir, por los enunciados que prescriben cómo deben ser las cosas. Por ejemplo, que está prohibido matar y que quien lo haga tendrá determinadas consecuencias, salvo que la propia norma prevea una excepción, como la legítima defensa.
\end{itemize}

La calificación jurídica de un hecho —por ejemplo, si una muerte constituye homicidio o legítima defensa— se produce en el momento en que ese hecho logra \textbf{subsumirse}, es decir, adecuarse, a la norma.

\section{La subsunción: del hecho a la norma}

Ante un hecho, el jurista debe mostrar y lograr convencer al juez de que existió un daño, un perjuicio, un cambio de estado: que una persona o un objeto se encontraba de determinada manera y, luego de un hecho, dejó de estarlo. Es necesario, además, identificar la causa que produjo esa modificación.

\begin{definition}{Daño}
Afectación de un derecho: modificación perjudicial de un estado de cosas, producida por un hecho con el que debe existir relación de causalidad.
\end{definition}

Si el daño fue producido por un fenómeno sin autor identificable —por ejemplo, un rayo—, resulta muy difícil encontrar un responsable. En cambio, si existió una conducta identificable —como la maniobra de una lancha— que provocó el daño, existe una relación de causalidad y, por lo tanto, un responsable. Lo primero que debe encontrarse, entonces, es que algo se hallaba de una manera y se modificó; esa modificación pertenece a alguno de los dos mundos ontológicos —el ser y el deber ser— con los que trabaja siempre el jurista.

\section{El caso base: el accidente de lancha en el Tigre}

A lo largo de la unidad se trabaja un caso guía para ejemplificar la construcción de una demanda: un accidente ocurrido entre una lancha comercial de pasajeros y una pequeña embarcación en el río Tigre. Existe un estado de cosas que, como consecuencia de un hecho de mayor o menor violencia, sufre una alteración; en el caso analizado, esa alteración consistió en una muerte y, además, en un daño material, aunque este último no resulte comparable con la pérdida de una vida humana.

\begin{important}
El razonamiento jurídico frente a un caso concreto consiste, en general, en verificar si un hecho produjo un daño subsumible en el principio general de no dañar a otro (\textit{alterum non laedere}), consagrado como fundamento de la responsabilidad civil en el Código Civil y Comercial de la Nación.
\end{important}

El código procesal estructura el contenido de la demanda siguiendo esta misma lógica: primero se ubica el mundo del ser —probar la realidad, los hechos efectivamente ocurridos— y después se ubica lo que se pide, la pretensión, correspondiente al mundo del deber ser. Esa estructura apunta, además, a que exista una organización que permita reclamar y que, en igualdad de condiciones, permita a la contraparte defenderse: debe existir «un espejo» entre la demanda y la contestación, en un tipo de discurso emparentado con la retórica clásica de la tradición griega y romana, cuyo objetivo consiste en estructurar el planteo del conflicto.

\section*{Cuadro Cornell — Fundamentos del proceso judicial}
\addcontentsline{toc}{section}{Cuadro Cornell — Fundamentos del proceso judicial}

\begin{table}[H]
\centering
\begin{tabularx}{\textwidth}{
    >{\raggedright\arraybackslash}p{0.28\textwidth}
    >{\raggedright\arraybackslash}X
}
\toprule
\textbf{Preguntas / palabras clave} &
\textbf{Notas / respuestas} \\
\midrule

\textbf{¿Qué es la jurisdicción?} &
Es la potestad exclusiva del juez de decir el derecho aplicable a un caso y de aplicarlo, incluso haciéndolo cumplir por la fuerza si la parte condenada no lo hace espontáneamente. \\

\textbf{¿Cuál es el objetivo del proceso?} &
Llegar a una sentencia que resuelva el conflicto, garantizando la igualdad entre las partes ante un tercero imparcial, ya que la justicia como ideal absoluto resulta un objetivo inalcanzable. \\

\textbf{¿Qué es el tema decidendum?} &
El límite dentro del cual puede resolver el juez, fijado por la pretensión de la parte actora, la defensa del demandado y la prueba producida en el proceso. \\

\textbf{¿Cómo se relacionan el «ser» y el «deber ser»?} &
El «ser» son los hechos ocurridos en la realidad; el «deber ser» es la norma jurídica. La tarea del operador jurídico consiste en subsumir el hecho en la norma para darle relevancia jurídica. \\

\textbf{¿Qué es el daño?} &
La afectación de un derecho: una modificación perjudicial de un estado de cosas, producida por un hecho con el que debe existir relación de causalidad. \\

\midrule

\multicolumn{2}{p{\dimexpr\textwidth-2\tabcolsep\relax}}{
\textbf{Resumen:} El proceso judicial busca garantizar la igualdad entre las partes frente a un tercero imparcial, porque la justicia como ideal absoluto es inalcanzable. El rol del abogado litigante es tan decisivo como el del juez, porque es quien construye el planteo —a través del hecho, la norma y la subsunción— que hace posible la sentencia, dentro del límite fijado por el tema decidendum.
} \\

\bottomrule
\end{tabularx}
\end{table}

% =========================================================================
% CAPÍTULO 2
% =========================================================================
\chapter[Retórica y estrategia en el litigio]{Retórica y estrategia en el litigio}

\section{El discurso jurídico como discurso persuasivo}

La demanda, y los juicios en general, constituyen una manifestación de lo que la tradición griega denominaba \textit{el discurso}, cuya estructura fija un contenido mínimo sin impedir agregar otras cuestiones. Preguntarse y responderse no es, en sí mismo, la retórica, sino un recurso propio de ella: la \textbf{pregunta retórica}.

Todo discurso persuasivo busca convencer a un destinatario. En el discurso político, el destinatario es el electorado; en el discurso jurídico, el destinatario es el juez. La diferencia entre ambos radica precisamente en la persona a la cual se dirige el discurso. Sin embargo, el deseo de convencer a otro no es exclusivo de políticos ni de abogados: es propio de la naturaleza humana y de la vida en sociedad, ya que cada vez que se solicita algo a otra persona —con mayor o menor honestidad— se le está ofreciendo una razón, construida a través de una situación retórica.

\begin{note}
El apunte ilustra esta idea con una anécdota: en una ciudad bilingüe, personas que pedían dinero en la calle repetían su pedido primero en un idioma y, ante la falta de respuesta, insistían en el otro. Más allá de lo anecdótico, el ejemplo evidencia que la estructura del discurso retórico —persuadir para obtener algo— es la misma que emplea cualquier persona que pide algo, dentro o fuera del ámbito jurídico.
\end{note}

\section{Lógica y emoción en la decisión judicial}

El discurso jurídico procura llegar al convencimiento por dos vías: una \textbf{lógica}, vinculada con la aplicación de las normas, y otra \textbf{emocional}, en tanto los seres humanos —incluidos los jueces— se relacionan y evalúan también a través de las emociones. Si bien existen juicios en los que la cuestión es obvia y no hay demasiado para discutir, en la enorme mayoría de los casos lo ideológico y lo emocional está presente durante todo el proceso de decisión.

\section{Caso: la revinculación paterno-filial y las preguntas retóricas}

\begin{example}{Revinculación paterno-filial}
Un hombre fue denunciado en Francia por presunto abuso contra su hijo varón. Tanto la causa penal como la civil tramitadas en Francia resultaron, con el tiempo, favorables a él. La justicia argentina había dispuesto que los niños regresaran al país, por considerar que allí estaba fijado su centro de vida. Al llegar los niños, la madre formuló una nueva denuncia, idéntica a la anterior, y a partir de entonces le retiraron la custodia al padre y se la otorgaron a la madre.

Desde ese momento —cinco años antes del relato— el padre no pudo ver ni una fotografía de sus hijos. Respecto de la niña se dictó sobreseimiento total y firme, sin que la contraparte lo hubiera apelado; respecto del hijo varón se dispuso la elevación a juicio, etapa en la que el proceso continuaba al momento del relato.
\end{example}

Se inició un incidente de revinculación, ofreciendo todas las garantías que la contraparte pudiera exigir —custodia policial, asistentes sociales, acompañamiento psicológico— con tal de poder comenzar algún grado de contacto. La contraparte sostuvo que el padre no debía ver a sus hijos, postura a la que adhirió el Ministerio Público. El juez resolvió en el mismo sentido, invocando el interés superior del niño, calificando lo solicitado como una medida cautelar y sosteniendo que no era el momento ni estaban dadas las condiciones para modificar la situación existente. En consecuencia, rechazó la revinculación.

\begin{example}{Pregunta retórica construida frente al caso}
«Señor juez, ¿dónde está escrito que el interés superior del niño consista en no ver nunca más a su papá, ni a su familia paterna, ni a sus tíos, ni a sus primos, ni a su abuela? ¿Dónde está escrito que usted, de esta manera, está ayudando al niño? Esto no es una sentencia: esto es, poco menos, una falta de respeto hacia el menor. ¿Dónde encontró usted en la ley que el papá no pueda ver una fotografía de su hijo, o comunicarse con él aunque sea por un medio audiovisual, con todos los recursos tecnológicos que existen hoy? ¿Cuál es el riesgo que usted encuentra en que el papá pueda ver a la nena a través de una pantalla?»
\end{example}

\begin{definition}{Pregunta retórica}
Recurso del discurso persuasivo mediante el cual quien lo formula se pregunta y se responde a sí mismo, con el fin de reforzar y responder aquello que está afirmando. Su función es enfatizar y dejar al interlocutor «totalmente en falta», sin posibilidad de disimular la arbitrariedad de una decisión.
\end{definition}

\section*{Cuadro Cornell — Retórica y estrategia en el litigio}
\addcontentsline{toc}{section}{Cuadro Cornell — Retórica y estrategia en el litigio}

\begin{table}[H]
\centering
\begin{tabularx}{\textwidth}{
    >{\raggedright\arraybackslash}p{0.28\textwidth}
    >{\raggedright\arraybackslash}X
}
\toprule
\textbf{Preguntas / palabras clave} &
\textbf{Notas / respuestas} \\
\midrule

\textbf{¿Para qué sirven las preguntas retóricas?} &
Son un recurso del discurso persuasivo que refuerza una idea y expone la arbitrariedad del contrario, sin que este pueda evadir la contradicción planteada. \\

\midrule

\multicolumn{2}{p{\dimexpr\textwidth-2\tabcolsep\relax}}{
\textbf{Resumen:} Convencer al juez exige tanto solidez lógica como sensibilidad hacia la dimensión emocional e ideológica que atraviesa toda decisión judicial, algo que el caso de la revinculación paterno-filial ilustra con crudeza a través del uso de preguntas retóricas frente a una decisión arbitraria.
} \\

\bottomrule
\end{tabularx}
\end{table}

% =========================================================================
% CAPÍTULO 3
% =========================================================================
\chapter[Estructura de la demanda]{Estructura de la demanda}

\section{El sumario y el encabezamiento}

La estructura de la demanda comienza por el sumario y el encabezamiento, que admiten distintas fórmulas equivalentes —«promueve demanda», «inicia demanda», «demanda»—, entre las cuales cada profesional puede optar. Dentro de esa estructura, el único dato que variará de un caso a otro es el domicilio real de la parte actora.

\section{El objeto de la demanda}

Después del encabezamiento, la demanda se dirige hacia el \textbf{objeto}, que tiene una importancia trascendente porque es exactamente lo que se pide. El objeto fija el tema decidendum desde la perspectiva de lo que la parte actora reclama: el juez no podrá conceder más de lo que allí se le pidió. Por esa razón, el objeto debe redactarse con claridad, precisión y concisión.

En el caso guía trabajado en la unidad, el objeto consiste en una indemnización por los daños sufridos: por la muerte de la señora accidentada y por los daños causados a su familia.

\begin{formula}{Objeto de la demanda}
«[Nombre de la parte actora] viene a promover demanda de daños y perjuicios contra [demandado/s], por la suma de [monto], en razón de los hechos y el derecho que se expondrán, como consecuencia del accidente relatado.»
\end{formula}

\section{La legitimación activa y pasiva}

\begin{definition}{Legitimación}
Condición de ser titular del derecho que se reclama (legitimación activa) o de ser la persona a quien puede exigírsele el cumplimiento de una obligación (legitimación pasiva).
\end{definition}

La legitimación se distingue en dos planos: la \textbf{legitimación activa} —quién puede reclamar el derecho en cuestión— y la \textbf{legitimación pasiva} —a quién puede dirigirse el reclamo—, ya que no es posible reclamar un derecho a cualquier persona.

\section{La elección de los demandados: criterios prácticos}

En el caso guía, los demandados identificados son los dueños de la lancha, la empresa naviera y la aseguradora; se descarta demandar a la municipalidad, aunque se analiza la posibilidad de sumar al Estado —por ejemplo, a la Prefectura Naval— como codemandado.

La cantidad de demandados no es indiferente: cada demandado implica dirigir una cédula de notificación distinta, mayor carga procesal y más tiempo. Por ello conviene evaluar previamente:

\begin{itemize}
    \item la solvencia de cada potencial demandado, es decir, quién está en mejores condiciones de pagar;
    \item la carga procesal que implica sumar cada codemandado;
    \item los plazos procesales aplicables —los que rigen contra el Estado, por ejemplo, son más largos—;
    \item la probabilidad de que ese demandado termine siendo responsable y, en definitiva, quien efectivamente pague.
\end{itemize}

Respecto de organismos estatales, puede convenir incluirlos aun cuando los plazos procesales contra el Estado sean más largos, porque en definitiva pueden terminar siendo responsables y, por lo tanto, quienes efectivamente paguen. Se trata, en suma, de una decisión estratégica que cada profesional debe evaluar en cada caso concreto.

\section{La narración de los hechos: redacción del caso práctico}

\begin{example}{Relato de hechos del caso guía (redacción realizada en clase)}
«El hecho que dio origen a la presente demanda ocurrió el 9 de [mes] de 2020, a las 11:30 horas aproximadamente, mientras las señoras Lorena Montes y Carla González circulaban a bordo de una pequeña embarcación por el río Tigre, perteneciente al Club Delta, cuando, cerca de la estación fluvial de Tigre, fueron embestidas por la lancha comercial de pasajeros de la línea Tigre Argentina S.R.L.»

«La lancha colectiva se encontraba realizando maniobras de marcha atrás para salir de la terminal y, al retroceder, succionó y pasó por encima de la pequeña embarcación. Como consecuencia del impacto, la señora Lorena Montes, luego de permanecer hospitalizada durante cinco días en un centro de salud de la localidad de Pacheco, falleció; la señora Carla González sufrió lesiones leves.»

«Los pormenores de la mecánica del siniestro surgirán con exactitud de las pruebas a producirse en autos, en especial de la pericia accidentológica. Lo sucedido dio origen a la causa penal caratulada «N.N. s/ homicidio culposo», que tramita ante el juzgado de instrucción en lo criminal y correccional interviniente.»
\end{example}

% TODO: contenido incompleto o ambiguo en las notas originales — el relato de hechos fue calificado por el propio docente como incompleto («es corta, porque son muchos hechos»); restan definir, entre otras cuestiones, aspectos relevantes como el lugar exacto de presentación de la demanda.

\section*{Cuadro Cornell — Estructura de la demanda}
\addcontentsline{toc}{section}{Cuadro Cornell — Estructura de la demanda}

\begin{table}[H]
\centering
\begin{tabularx}{\textwidth}{
    >{\raggedright\arraybackslash}p{0.28\textwidth}
    >{\raggedright\arraybackslash}X
}
\toprule
\textbf{Preguntas / palabras clave} &
\textbf{Notas / respuestas} \\
\midrule

\textbf{¿Qué exige el artículo 330 del CPCCN?} &
Fija el contenido mínimo que debe tener toda demanda: encabezamiento, hechos, derecho, petitorio (objeto) y ofrecimiento de prueba. \\

\textbf{¿Qué es la legitimación?} &
La condición de ser titular del derecho que se reclama (legitimación activa) o de ser la persona a quien puede exigírsele el cumplimiento de una obligación (legitimación pasiva). \\

\textbf{¿Qué criterios definen a quién demandar?} &
La solvencia del demandado, los plazos procesales aplicables, la carga procesal que implica cada codemandado y la probabilidad de éxito frente a cada uno. \\

\midrule

\multicolumn{2}{p{\dimexpr\textwidth-2\tabcolsep\relax}}{
\textbf{Resumen:} Sobre la base de la lógica hecho-norma expuesta en el capítulo anterior, se construye la estructura formal de la demanda —sumario, objeto, legitimación y hechos— tal como la exige el artículo 330 del CPCCN, aplicada de manera práctica al caso guía del accidente de lancha en el río Tigre.
} \\

\bottomrule
\end{tabularx}
\end{table}

% =========================================================================
% CAPÍTULO 4
% =========================================================================
\chapter[Competencia y organización de la justicia argentina]{Competencia y organización de la justicia argentina}

\section{La competencia territorial: los distritos judiciales}

En el caso guía, las aseguradoras tienen su domicilio en la Capital, lo cual habilita una primera gran decisión: si demandar ante la justicia nacional o ante la provincial. Tratándose de un hecho ocurrido en Tigre, la demanda podría iniciarse ante la justicia provincial correspondiente.

La Provincia de Buenos Aires está organizada en distintos distritos judiciales, que no necesariamente coinciden con los partidos o municipios. El distrito judicial de San Isidro, por ejemplo, comprende los partidos de Tigre, Vicente López, San Isidro y San Fernando; el de Lomas de Zamora comprende ese partido y otros vecinos, y así sucesivamente. En toda la provincia se aplica el mismo Código Procesal, pero los tribunales tienen asiento físico en un municipio determinado dentro de cada distrito. En el caso trabajado, tratándose de un hecho ocurrido en Tigre, la demanda podría iniciarse ante los tribunales de San Isidro.

\section{Justicia ordinaria y justicia federal: el reparto constitucional}

En todo el territorio del país coexisten siempre dos órdenes superpuestos: el federal y el ordinario; no hay zonas del país en las que exista uno sin el otro. La justicia federal es de aplicación excepcional —es decir, extraordinaria—, lo cual supone que la mayoría de las causas son ordinarias.

\begin{important}
Los artículos 116 y 117 de la Constitución Nacional determinan cuáles son las cuestiones que corresponden a la competencia de la justicia federal. A modo de ejemplo: un conflicto entre dos provincias corresponde a la Corte Suprema de Justicia de la Nación —de manera originaria o por apelación, según el caso—; una cuestión que involucre a un cónsul o a un embajador corresponde a la justicia federal; y lo mismo ocurre, en principio, con cuestiones de telecomunicaciones o de ferrocarriles.
\end{important}

La justicia ordinaria depende de las provincias, por tratarse de una facultad no delegada al gobierno federal. En algún momento las distintas provincias fueron, de hecho, estados independientes entre sí —incluso Buenos Aires funcionó como estado propio antes de 1860—. Cuando esas provincias decidieron constituirse en una Nación, lo hicieron a través de un pacto, y en ese pacto delegaron ciertos poderes al gobierno federal —como la creación de leyes de fondo, la organización de la justicia federal y la administración a cargo del poder ejecutivo nacional— pero se reservaron otras facultades, entre ellas, la organización de su propia administración de justicia. Por eso cada provincia cuenta con su propia Constitución y con sus propios tres poderes —ejecutivo, legislativo y judicial— y con un proceso propio, distinto del de las demás provincias, aunque compartan buena parte de su estructura: primera instancia, cámara de apelaciones y, eventualmente, una Corte o Superior Tribunal de Justicia provincial.

Las cuestiones ordinarias surgen del Código Civil y Comercial, del Código Penal, y también de leyes especiales de aplicación uniforme en todo el país —como la Ley de Contrato de Trabajo—, que, aunque no son formalmente un código, funcionan como tal: la ley de fondo es la misma en toda la Nación, pero se aplica a través de procesos diferentes según cada jurisdicción. Frente a supuestos dudosos —como un choque ocurrido en una ruta nacional—, la respuesta no depende de contra quién se litiga, sino de qué norma resulta aplicable al caso; en general, ese tipo de supuestos tramita ante la justicia ordinaria, aunque la cuestión admite discusión.

\section{La ley aplicable: código civil o ley de navegación}

En el caso guía, además del Código Civil, existe otra norma en discusión: la Ley de Navegación, en tanto pueda calificarse el hecho como un \textbf{abordaje}. No se trata de un choque de esquina: si la lancha hubiese pasado por encima de la plataforma sin llegar a producirse un contacto entre los cascos, ya no se plantearía el mismo problema de calificación.

\begin{important}
Existe una discusión doctrinaria y jurisprudencial acerca de si corresponde calificar como abordaje —lo que atraería la aplicación de la Ley de Navegación— un siniestro en el que intervienen una lancha comercial de gran porte y una embarcación pequeña. Esa calificación incide, a su vez, en la competencia: puede atraer la intervención de la justicia federal.
\end{important}

El solo hecho de interponer la demanda ante un fuero determinado implica proponer esa jurisdicción, fijándola, aunque después la calificación jurídica del hecho no se confirme. En el caso guía, las posibilidades de jurisdicción quedan resumidas en el siguiente cuadro:

\begin{table}[H]
\centering
\begin{tabularx}{\textwidth}{
    >{\raggedright\arraybackslash}p{0.55\textwidth}
    >{\raggedright\arraybackslash}X
}
\toprule
\textbf{Opción de jurisdicción} & \textbf{Fuero y asiento} \\
\midrule
1. Justicia provincial ordinaria & Civil y comercial ordinaria — San Isidro (Provincia de Buenos Aires) \\
2. Justicia federal con asiento en la Provincia & Fuero federal — San Martín \\
3. Justicia nacional ordinaria & Civil y comercial — Capital Federal \\
4. Justicia federal con asiento en la Capital & Fuero federal — Capital Federal \\
\bottomrule
\end{tabularx}
\end{table}

\section*{Cuadro Cornell — Competencia y organización de la justicia argentina}
\addcontentsline{toc}{section}{Cuadro Cornell — Competencia y organización de la justicia argentina}

\begin{table}[H]
\centering
\begin{tabularx}{\textwidth}{
    >{\raggedright\arraybackslash}p{0.28\textwidth}
    >{\raggedright\arraybackslash}X
}
\toprule
\textbf{Preguntas / palabras clave} &
\textbf{Notas / respuestas} \\
\midrule

\textbf{¿Cómo se determina la competencia territorial?} &
Por el lugar del hecho o por el domicilio del demandado, dentro de la organización de distritos judiciales de cada provincia. \\

\textbf{¿Qué diferencia a la justicia ordinaria de la federal?} &
La ordinaria es la regla general y depende de las provincias, por tratarse de una facultad no delegada (arts. 116 y 117 CN); la federal es de aplicación excepcional y se limita a las materias y personas reservadas a la Nación. \\

\midrule

\multicolumn{2}{p{\dimexpr\textwidth-2\tabcolsep\relax}}{
\textbf{Resumen:} La unidad desarrolla el mapa de competencias de la Argentina —la distinción entre justicia ordinaria y justicia federal, fundada en el reparto constitucional de facultades no delegadas por las provincias— y su aplicación práctica al caso guía, donde la calificación del hecho como abordaje frente a la Ley de Navegación condiciona el fuero ante el que corresponde demandar.
} \\

\bottomrule
\end{tabularx}
\end{table}

% =========================================================================
% CAPÍTULO 5
% =========================================================================
\chapter[La justicia de la Ciudad Autónoma de Buenos Aires]{La justicia de la Ciudad Autónoma de Buenos Aires}

Además de la distinción entre justicia ordinaria y justicia federal, la organización judicial de la Ciudad de Buenos Aires presenta una complejidad adicional: el juez civil de la Capital termina siendo pagado con los impuestos de todo el país, y no solamente por los habitantes de la Ciudad.

\section{Origen histórico: la Ley Avellaneda de 1880}

\begin{note}
Hacia 1880, mediante la denominada Ley Avellaneda, se federalizó la ciudad de Buenos Aires. Hasta entonces, las primeras autoridades nacionales tenían asiento en un territorio que formaba parte de la Provincia de Buenos Aires, lo cual generaba una superposición de poderes. Para resolverlo, la provincia cedió el territorio de la ciudad, que pasó a constituir la Capital Federal, y se creó la ciudad de La Plata como nueva capital de la Provincia de Buenos Aires.
\end{note}

A partir de entonces, en la Ciudad de Buenos Aires existieron jueces «nacionales» —distintos de los jueces federales— encargados de resolver las cuestiones ordinarias de sus habitantes, ya que la Ciudad, al no ser una provincia, carecía de un gobierno propio que organizara su propia justicia.

\section{La autonomía porteña y la creación de nuevos tribunales}

Antes de 1994, la Ciudad de Buenos Aires se gobernaba mediante una intendencia cuyo titular era designado por el presidente de la Nación. Con la reforma constitucional de 1994, la Ciudad adquirió el carácter de Ciudad Autónoma, equiparable en varios aspectos a una provincia, con Constitución propia, poder legislativo propio —diputados y senadores— y representación de tres senadores, igual que el resto de las provincias.

Al equipararse a una provincia, correspondía que la Ciudad tuviera su propia justicia. Los jueces nacionales que ya ejercían en la Ciudad se resistieron a integrar una justicia local, por lo que se optó por crear tribunales nuevos: los tribunales en lo Contencioso Administrativo y Tributario de la Ciudad, con sus respectivas Cámaras de apelación y, en la cúspide, el Tribunal Superior de Justicia porteño.

\begin{example}{Distinción entre fuero ordinario y contencioso administrativo porteño}
Si una persona se atiende en un sanatorio privado —como el Hospital Alemán— y sufre una mala praxis, debe demandar ante la justicia ordinaria. Si, en cambio, se atiende en un hospital público de la Ciudad —como el Hospital Fernández— y ocurre lo mismo, debe demandar ante el fuero contencioso administrativo y tributario porteño, porque en ese caso la propia Ciudad es parte demandada.
\end{example}

\section{El fallo «Vinas» y el fuero contencioso administrativo}

En el resto de las provincias, un proceso avanza de primera instancia a la Cámara de Apelaciones y luego, eventualmente, al Tribunal Superior o Corte provincial, antes de poder llegar a la Corte Suprema de Justicia de la Nación. En la Capital Federal, en cambio, los procesos que salían de los tribunales nacionales iban directamente a la Corte Suprema.

\begin{important}
El fallo «Vinas» de la Corte Suprema de Justicia de la Nación dispuso que, antes de llegar a la Corte, los procesos originados en los tribunales nacionales de la Ciudad deben pasar primero por el Tribunal Superior de Justicia de la Ciudad de Buenos Aires, generando así una instancia intermedia adicional.
\end{important}

\section{La reforma laboral y el traspaso de competencias a partir de 2027}

A partir de 2027, con la nueva ley, los juicios laborales que se inicien tramitarán ante los nuevos tribunales creados en la justicia de la Ciudad Autónoma de Buenos Aires. El fuero laboral nacional continuará entendiendo únicamente en las causas iniciadas hasta el 31 de diciembre de 2026; los juicios laborales iniciados a partir del 1 de enero de 2027 tramitarán ante la justicia porteña.

% TODO: contenido incompleto o ambiguo en las notas originales — el apunte no precisa el número o la denominación oficial de la ley que dispone este traspaso de competencias, ni el detalle completo de su régimen transitorio.

\section*{Cuadro Cornell — La justicia de la Ciudad Autónoma de Buenos Aires}
\addcontentsline{toc}{section}{Cuadro Cornell — La justicia de la Ciudad Autónoma de Buenos Aires}

\begin{table}[H]
\centering
\begin{tabularx}{\textwidth}{
    >{\raggedright\arraybackslash}p{0.28\textwidth}
    >{\raggedright\arraybackslash}X
}
\toprule
\textbf{Preguntas / palabras clave} &
\textbf{Notas / respuestas} \\
\midrule

\textbf{¿Por qué existe una justicia «nacional» distinta de la «federal» en la Ciudad de Buenos Aires?} &
Por el proceso histórico de federalización de 1880 (Ley Avellaneda) y la posterior autonomía porteña de 1994, que generó una superposición de fueros nacionales, federales y locales todavía no resuelta por completo. \\

\textbf{¿Qué estableció el fallo «Vinas»?} &
Que los procesos que salen de los tribunales nacionales de la Ciudad deben pasar primero por el Tribunal Superior de Justicia porteño, antes de poder llegar, eventualmente, a la Corte Suprema. \\

\midrule

\multicolumn{2}{p{\dimexpr\textwidth-2\tabcolsep\relax}}{
\textbf{Resumen:} La unidad cierra con el caso particular —históricamente complejo— de la organización judicial de la Ciudad Autónoma de Buenos Aires, desde la Ley Avellaneda de 1880 hasta el fallo «Vinas» y la reforma laboral que comenzará a regir en 2027, lo que confirma que redactar una demanda no es completar un formulario, sino tomar decisiones estratégicas fundadas —qué se pide, a quién y ante qué juez— que definen, desde el primer escrito, el resultado de todo el proceso.
} \\

\bottomrule
\end{tabularx}
\end{table}

% =========================================================================
% CAPÍTULO 6
% =========================================================================
\chapter[Trabajo práctico integrador]{Trabajo práctico integrador}

Como cierre de la unidad, se propone un trabajo práctico integrador consistente en plantear y redactar la demanda del caso guía trabajado a lo largo de la clase.

\begin{exercise}{Consigna para la clase siguiente}
A partir del caso guía del accidente de lancha en el río Tigre, se solicita:

\begin{enumerate}
    \item Plantear la competencia del caso, fundamentando la elección del fuero.
    \item Definir a quiénes se va a demandar.
    \item Redactar la narración completa de los hechos, de forma clara, ordenada y cronológica.
    \item Elaborar el objeto de la demanda conforme a la fórmula trabajada en la unidad.
\end{enumerate}
\end{exercise}

\begin{note}
La unidad siguiente completará la estructura de la demanda con el desarrollo de la prueba: qué más lleva la demanda, además de los hechos, y un repaso de cada uno de los medios de prueba.
\end{note}

\end{document}
